\documentclass[aspectratio=169,12pt]{beamer}

\usepackage{fontspec}
\usepackage{polyglossia}
\setdefaultlanguage{russian}
\setmainfont{PTSans.ttc}[Path=/System/Library/Fonts/Supplemental/]
\setsansfont{PTSans.ttc}[Path=/System/Library/Fonts/Supplemental/]
\setmonofont{PTMono.ttc}[Path=/System/Library/Fonts/Supplemental/]

\usepackage{booktabs}
\usepackage{array}
\usepackage{ragged2e}

\definecolor{vtbblue}{HTML}{0F2747}
\definecolor{vtborange}{HTML}{FF8C42}
\definecolor{vtblight}{HTML}{F6F1E8}
\definecolor{vtbtext}{HTML}{10253F}

\usetheme{Madrid}
\setbeamercolor{normal text}{fg=vtbtext,bg=white}
\setbeamercolor{structure}{fg=vtbblue}
\setbeamercolor{title}{fg=white,bg=vtbblue}
\setbeamercolor{frametitle}{fg=white,bg=vtbblue}
\setbeamercolor{block title}{fg=white,bg=vtbblue}
\setbeamercolor{block body}{fg=vtbtext,bg=vtblight}
\setbeamercolor{item projected}{fg=white,bg=vtborange}
\setbeamertemplate{navigation symbols}{}
\setbeamertemplate{itemize item}{\color{vtborange}$\blacksquare$}
\setbeamertemplate{itemize subitem}{\color{vtborange}$\bullet$}

\title{Выбор новых зон для банкоматов ВТБ в Москве}
\subtitle{Баланс плотности сети, экономического эффекта и сервисного покрытия}
\author{VTB Geo Analytics}
\date{Май 2026}

\begin{document}

\begin{frame}
  \titlepage
\end{frame}

\begin{frame}[shrink=10]{Как решалась задача баланса}
\footnotesize
\begin{columns}[T]
\begin{column}{0.60\textwidth}
\begin{block}{Постановка задачи}
\begin{itemize}
\item Нельзя брать только локально сильные точки: центр быстро перегревается.
\item Новые устройства начинают конкурировать друг с другом и с действующей сетью ВТБ.
\item Поэтому задача решалась как \textbf{последовательное расширение сети}, а не как выбор фиксированного набора \texttt{top-K}.
\end{itemize}
\end{block}

\begin{block}{Методология}
\begin{enumerate}
\item Москва разбита на \textbf{8\,154 H3-ячейки}; для каждой ячейки рассчитаны спрос, покрытие и подходящий тип ATM.
\item На каждом шаге оптимизатор считает \textbf{предельную полезность} новой точки и вычитает штрафы за близость к сети ВТБ и чрезмерный рост сети.
\end{enumerate}
\end{block}
\end{column}

\begin{column}{0.40\textwidth}
\begin{block}{Результат по сценариям}
\centering
\begin{tabular}{>{\raggedright\arraybackslash}p{0.47\linewidth} >{\centering\arraybackslash}p{0.28\linewidth}}
\toprule
Сценарий & Оптимум, новых ATM \\
\midrule
Сбалансированный & \textbf{65} \\
Прибыль & 76 \\
Покрытие & 50 \\
Социальный & 40 \\
Конкурентный & 56 \\
Малый бизнес & 56 \\
\bottomrule
\end{tabular}
\end{block}

\begin{alertblock}{Ключевой вывод}
В сбалансированном режиме максимум функции полезности достигается после \textbf{+65 новых банкоматов}. Следовательно, сеть ВТБ в Москве нужно расширять.
\end{alertblock}
\end{column}
\end{columns}
\end{frame}

\begin{frame}[shrink=8]{Данные и модельный конвейер}
\footnotesize
\begin{block}{Как работали с данными}
\begin{itemize}
\item Внутренний слой: \textbf{4.15 млн строк транзакций}, \textbf{69\,337 клиентов}, агрегация до H3.
\item Целевой слой: ATM-активность и интенсивность спроса в ячейках.
\item Внешнее обогащение: OpenStreetMap, метро, школы, вузы, больницы, торговля, офисы, плотность чужих ATM, близость к присутствию ВТБ.
\item Итоговая витрина: \textbf{131 признак} на \textbf{8\,154 ячейках}.
\end{itemize}
\end{block}

\begin{block}{Модельный конвейер}
\begin{enumerate}
\item \textbf{CatBoostClassifier} оценивает вероятность сильной ATM-активности в ячейке.
\item \textbf{CatBoostRegressor} прогнозирует интенсивность ATM-спроса.
\item Оба прогноза объединяются в \textbf{составной балл ячейки}.
\item Затем подбирается один из четырех архетипов ATM: стандартный, полнофункциональный, high-flow или business cash-in.
\item Поверх баллов работает \textbf{сетевой оптимизатор}, который выбирает лучший набор новых точек.
\end{enumerate}
\end{block}
\end{frame}

\begin{frame}{Почему выбрана именно такая архитектура}
\footnotesize
\begin{columns}[T]
\begin{column}{0.52\textwidth}
\begin{block}{Архитектурные причины}
\begin{itemize}
\item Данные табличные и разнородные: градиентный бустинг здесь надежнее тяжелой геомодели.
\item Пространственная структура не теряется: она зашита в признаки через H3-соседство, расстояния, плотности и пространственные лаги.
\item Архитектура интерпретируема: можно объяснить вклад признаков и логику выбора точки.
\item Такая схема удобна для защиты: отдельно объясняются спрос, покрытие, тип ATM и итоговый сетевой выбор.
\end{itemize}
\end{block}
\end{column}
\begin{column}{0.48\textwidth}
\begin{block}{Бизнес-режимы}
\begin{itemize}
\item \textbf{Сбалансированный}: доход + покрытие + контроль плотности.
\item \textbf{Прибыль}: максимизация денежного потенциала.
\item \textbf{Покрытие}: закрытие сервисных пустот.
\item \textbf{Социальный}: доступность рядом со школами и медициной.
\item \textbf{Конкурентный}: ответ на давление чужих сетей.
\item \textbf{Малый бизнес}: приоритет cash-in и торгово-офисных кластеров.
\end{itemize}
\end{block}
\end{column}
\end{columns}
\end{frame}

\begin{frame}[shrink=4]{Как подтверждали качество модели}
\footnotesize
\begin{columns}[T]
\begin{column}{0.54\textwidth}
\begin{block}{Схема проверки}
\begin{itemize}
\item Использована \textbf{пространственная кросс-проверка на 5 блоках}: обучение и проверка делятся по географическим зонам.
\item Для классификации считались \textbf{AUC}, средняя точность, \textbf{Precision@K}, \textbf{Recall@K}.
\item Для ранжирования и интенсивности спроса считались \textbf{RMSE}, \textbf{MAE}, \textbf{NDCG@K}.
\item Модель сравнивалась не только сама с собой по фолдам, но и с простым базовым уровнем.
\end{itemize}
\end{block}

\end{column}

\begin{column}{0.46\textwidth}
\begin{block}{Средние метрики по 5 пространственным блокам}
\centering
\begin{tabular}{>{\raggedright\arraybackslash}p{0.52\linewidth} >{\centering\arraybackslash}p{0.20\linewidth}}
\toprule
Показатель & Значение \\
\midrule
AUC & \textbf{0.910} \\
Средняя точность & \textbf{0.729} \\
Precision@K & \textbf{0.700} \\
Recall@K & \textbf{0.520} \\
RMSE спроса & \textbf{50.88} \\
NDCG@K & \textbf{0.782} \\
\bottomrule
\end{tabular}
\end{block}

\begin{block}{Сравнение с базовым уровнем}
\begin{itemize}
\item AUC: \textbf{0.910} против \textbf{0.859}
\item Средняя точность: \textbf{0.729} против \textbf{0.688}
\item RMSE спроса: \textbf{50.88} против \textbf{156.89}
\end{itemize}
\end{block}
\end{column}
\end{columns}
\end{frame}

\end{document}
